\documentclass[11pt]{article}
\usepackage[margin=0.9in]{geometry}
\usepackage{graphicx,amsmath,amssymb,booktabs,hyperref,xcolor}
\graphicspath{{/tmp/cheby_h0/figs/tau_report/}}
\hypersetup{colorlinks=true,linkcolor=blue,urlcolor=blue}
\setlength{\parskip}{6pt}
\title{Lin-CDF Richardson: 100-$\tau$ sweep at $\gamma=1$\\
       (with clipping-artifact warning at high $\tau$)}
\author{Fixed-Point Factory --- signal precision sweep}
\date{June 2026}

\begin{document}
\maketitle

\section*{Headline}

100 values of signal precision $\tau \in [0.01, 1000]$ at fixed
$\gamma = 1$, solved with Lin-CDF 2-pt Richardson and warm-start chain.
\textbf{99/100 reach machine $\varepsilon$.} Per-$\tau$ wall time $\sim 1$\,s.

\paragraph{Two regimes:}
\begin{itemize}
  \item \textbf{Valid range $\tau \lesssim 5$}: clean economics.
        Slope $\alpha^*$ rises from $0.33$ at $\tau{=}0.01$ to a peak
        of $0.44$ around $\tau{=}3$. One-to-one deficit grows from
        $\sim 0$ to $\sim 0.05$ over the same range --- partial revelation
        emerges as signals become more informative.
  \item \textbf{High-$\tau$ clipping artifact ($\tau \gtrsim 5$)}: slope
        drops toward 0, deficit flattens at $\sim 0.11$. \emph{Not}
        economically meaningful --- this is a numerical clipping artifact
        ($1\!-\!P$ clipped at $10^{-15}$ caps $|\mathrm{logit}\,P| \leq 34.5$).
        At high $\tau$ the true $\mathrm{logit}\,P$ at corner cells
        would be $\sim 0.5 \tau u_{\max}^2 \sim 100$+, so most cells
        saturate at the clip.
\end{itemize}

\section{Convergence}

\begin{center}
\includegraphics[width=\linewidth]{02_floor.png}
\end{center}

99/100 reach machine $\varepsilon$. The lone failure at $\tau \approx 3.4$
is a transitional regime (intermediate signal precision with mixed
revelation/non-revelation features).

\section{Equilibrium statistics across $\tau$}

\begin{center}
\includegraphics[width=\linewidth]{01_headline.png}
\end{center}

The grey shaded region marks $\tau \geq 5$ where the metrics become
unreliable due to logit clipping. Within the valid range $\tau < 5$,
both slope and deficit are monotone in $\tau$.

\section{Zoomed: clipping-free range $\tau < 5$}

\begin{center}
\includegraphics[width=\linewidth]{03_valid_range.png}
\end{center}

\textbf{Interpretation.}
\begin{itemize}
  \item As $\tau$ grows from $0$ to $\sim 3$, signals become more
        informative. The slope $\alpha^*$ of $\mathrm{logit}\,P$ on $T$
        rises because both the variance of $T$ and the price's
        response to it grow.
  \item The one-to-one deficit $1-R^2_{\rm nonparam}$ also rises ---
        partial revelation emerges in this regime. At very low $\tau$,
        signals are so noisy that the price is essentially $\sigma(T)$
        exactly (no partial revelation: $T$ contains all information).
        As $\tau$ grows, the individual $u_k$ start carrying information
        beyond what $T$ summarizes.
\end{itemize}

\newpage

\section{Linear vs nonparametric deficit}

\begin{center}
\includegraphics[width=0.95\linewidth]{05_deficits.png}
\end{center}

The linear deficit (vs $\alpha T$ fit) is always larger than the
nonparametric deficit (vs best $f(T)$). The gap is the residual
nonlinearity in $f(T)$ which has no economic meaning. The nonparametric
deficit is the right metric.

\section{logit-$P$ vs $T$ scatter (valid $\tau$ range)}

\begin{center}
\includegraphics[width=\linewidth]{06_logitT.png}
\end{center}

At low $\tau$ the scatter is tight (price is function of $T$). As $\tau$
grows the spread increases (different signal triples with same $T$
give different prices). Above $\tau \sim 5$ the linear regression
saturates due to clipping.

\section{Contour gallery: how the FP deforms with $\tau$}

\begin{center}
\includegraphics[width=\linewidth]{04_contours.png}
\end{center}

\section{Summary and recommendations}

\begin{center}
\begin{tabular}{lll}
\toprule
$\tau$ range & convergence & metric reliability \\
\midrule
$[0.01, 0.5]$ & machine $\varepsilon$ & reliable (clean economics) \\
$[0.5, 5]$    & machine $\varepsilon$ (except $\tau\approx 3.4$) & reliable \\
$[5, 1000]$   & machine $\varepsilon$ & \textbf{clipping artifact in slope/deficit} \\
\bottomrule
\end{tabular}
\end{center}

\paragraph{For the paper:}
\begin{itemize}
  \item Report $\tau \in [0.01, 5]$ economics.
  \item Show that slope $\alpha^*$ rises from $0.33$ to peak $\sim 0.44$
        across this range, and deficit rises from $0$ to $\sim 0.05$.
  \item Note that the operator continues to converge at high $\tau$ but
        the metrics are obscured by the float64 logit clip.
\end{itemize}

\paragraph{Comparison to $\gamma$-sweep at $\tau=1$:}
The $\tau$-sweep and $\gamma$-sweep both show monotonically rising slope
with the relevant parameter. But the slope direction is different:
\begin{itemize}
  \item Higher $\gamma$ (more risk averse) $\rightarrow$ price more revealing
        (higher slope). Mechanism: risk aversion damps demand response,
        the price is gentler.
  \item Higher $\tau$ (more precise signals) $\rightarrow$ same direction
        (higher slope), but for a different reason: agents' posteriors
        sharpen and the price reflects this.
\end{itemize}

\paragraph{Implementation note.} The Lin-CDF 2-pt Richardson operator
nailed 99 of 100 $\tau$ values in $\sim 100$ seconds total (warm-start
chain). This is the right tool for economic parametric sweeps.

\end{document}
